% ==========================================
%  content/02-abstract.tex  —  摘要与目录
%  可独立编译: xelatex 02-abstract.tex
% ==========================================
\documentclass[../main.tex]{subfiles}
\begin{document}

% --- 摘要页设置 ---
\renewcommand{\baselinestretch}{1.5}
\pagenumbering{Roman}
\pagestyle{fancy}
\fancyhead{}
\fancyhead[CE]{\songti{\zihao{5} \youreduction 士学位论文}}
\fancyhead[CO]{\songti{\zihao{5} \yourtitlenamechinese}}
\renewcommand{\headrulewidth}{1.5pt}
\renewcommand{\footrulewidth}{0pt}

% ==========================================
%  中文摘要
% ==========================================
\phantomsection\addcontentsline{toc}{chapter}{\heiti \zihao{4} 摘 \hspace*{1em} 要}\tolerance=500

\begin{cnabstract}
	{
		% 在此处填写中文摘要正文
		引力波的发现验证了广义相对论对黑洞与引力波的两大预言，开启了多信使天文学新时代。
		致密双星是最典型的引力波源，LIGO-Virgo-KAGRA（LVK）合作组织已确认超过100个致密双星并合事件。
		双星旋进阶段可通过后牛顿近似、后闵可夫斯基近似、引力自力等方法近似求解爱因斯坦场方程建模；
		旋进后期、并合与铃宕阶段则需采用数值相对论精确计算。
		因此，解析与数值方法需协同结合，这主要通过有效单体理论实现。

		\vspace{0.5cm}
		\par {\heiti \zihao{-4} 关键字: 后闵可夫斯基近似、有效单体理论、引力波、黑洞微扰、波形模板}
	}
\end{cnabstract}

\cleardoublepage

% ==========================================
%  英文摘要
% ==========================================
\phantomsection\addcontentsline{toc}{chapter}{\textbf{\zihao{4}ABSTRACT}}\tolerance=500

\begin{enabstract}
	{
		% 在此处填写英文摘要正文
		The discovery of gravitational waves has confirmed two major predictions of general relativity regarding
		black holes and gravitational waves, ushering in a new era of multi-messenger astronomy.
		Compact binary stars are the most typical sources of gravitational waves, and the LIGO-Virgo-KAGRA
		(LVK) collaboration has confirmed over 100 compact binary merger events.

		\vspace{0.5cm}
		\par {\bfseries {\zihao{-4} Keywords: Post-Minkowski Approximation, Effective-One-Body Theory,
			Gravitational Waves, Black Hole Perturbations, Waveform template}}
	}
\end{enabstract}

% ==========================================
%  目录
% ==========================================
\cleardoublepage
\begin{spacing}{2}
	\centering{\tableofcontents}
\end{spacing}

% 插图目录（如有需要取消注释）
% \cleardoublepage
% \listoffigures

% 表格目录（如有需要取消注释）
% \cleardoublepage
% \listoftables

\end{document}
